\section{Dispositions}
\label{dispositions}

\subsection{Principes directeurs}
\label{dispositions-principes-directeurs}

\begin{enumerate}
 \item Conformément aux statuts constitutifs, la FCEG soutient la croissance et la communication au sein de ses membres afin d'assurer leur bien-être moral, intellectuel, culturel, académique, social et économique.
 \item La FCEG représente et promeut la population étudiante de premier cycle en génie au Canada aux niveaux national et international.
 \item La FCEG ne s'alignera sur aucun parti politique.
\end{enumerate}

\subsection{Mission}
\label{dispositions-mission}

\begin{enumerate}
 \item L'objectif de la FCEG est de solliciter, de représenter, de défendre et de faire progresser les intérêts de ses membres au niveau national afin d'assurer leur bien-être moral, intellectuel, culturel, académique, social et économique, et d'améliorer la qualité et la portée de la formation en génie au Canada.
 \item La FCEG poursuivra cette mission en :
  \begin{enumerate}
   \item Organisant et en accueillant des congrès, des réunions, des assemblées, des expositions et des compétitions pour l'avancement de la mission,
   \item S'affiliant à toute organisation poursuivant des objectifs alignés sur ceux de la FCEG,
   \item Acquérant par achat, location ou autrement, en possédant et en exploitant les actifs, biens meubles et immeubles nécessaires à l'accomplissement de la mission,
   \item Fournissant à ses membres des services de toutes sortes liés à la mission,
   \item Créant, soutenant et participant à des mouvements, des campagnes d'opinion ou d'autres organisations faisant avancer ces objectifs,
   \item Encourageant les membres à créer et à partager des informations pertinentes avec tous les membres et autres parties intéressées,
   \item Promouvant l'interaction de la population étudiante en génie du Canada avec des groupes d'intérêt spécifiques (associations, gouvernements et autres agences) sur des questions nationales et internationales d'ordre social, économique, politique, juridique et humain pertinentes pour le génie,
   \item Présentant des informations pertinentes aux gouvernements et autres organismes appropriés,
   \item Coopérant avec toutes les associations de génie reconnues par l'organisation dans les questions d'intérêt commun,
   \item Promouvant et en sauvegardant l'image publique de la population étudiante en génie au Canada,
   \item Facilitant et en développant des activités et des services.
  \end{enumerate}
\end{enumerate}

\subsection{Régions}
\label{dispositions-regions}

\begin{enumerate}
 \item La FCEG est composée de quatre (4) régions, chacune représentée au niveau national :
  \begin{enumerate}
   \item Atlantique – Provinces de l’Île-du-Prince-Édouard, Nouveau-Brunswick, Nouvelle-Écosse et Terre-Neuve-et-Labrador;
   \item Québec – Province du Québec;
   \item Ontario – Province de l’Ontario; et
   \item Ouest – Provinces de Colombie-Britannique, Alberta, Saskatchewan, Manitoba, et le territoire du Yukon, le territoire du Nunavut et les Territoires du Nord-Ouest.
  \end{enumerate}
\end{enumerate}

\subsection{Langues officielles}
\label{dispositions-langues-officielles}

\begin{enumerate}
 \item La FCEG est une organisation bilingue. L'anglais et le français sont ses langues officielles.
\end{enumerate}
